% Section de traçabilité placée après les résultats.
% =============================================================================
\section{Traçabilité et archive scientifique}
\label{ann:tracabilite}
% =============================================================================

La livraison finale est regroupée dans le dépôt GitHub
\href{https://github.com/MaximDP/CS_Internship_Quantum_Paraelectrics}
{\nolinkurl{CS_Internship_Quantum_Paraelectrics}}. Celui-ci contient le rapport,
la soutenance Beamer, la chaîne de calcul SCHA effectivement exécutée, les jeux
de paramètres, les données sources disponibles, les figures livrées et les
instructions de reproduction. En l'absence de tag de livraison, un état
reproductible doit être identifié par son commit Git.

Les éléments nécessaires à la reproduction sont répartis comme suit :
\begin{itemize}
  \item \texttt{rapportCS\_stage\_maxime/biblio.bib} contient les références citées ;
  \item les dossiers \texttt{images/} et \texttt{figures/} du rapport et
  le dossier de figures du beamer contiennent les illustrations ;
  \item dans la chaîne de calcul, \texttt{inputs/} regroupe les configurations et données,
  et \texttt{outputs/reference/} conserve les résultats numériques livrés ;
  \item les chapitres théoriques du présent rapport exposent dans le corps du
  document les calculs, conventions et contrôles qui fondent les résultats.
  \item l'annexe~\ref{sec:core_parameter_sets} réunit les jeux de paramètres,
  les unités et l'ordre de reconstruction numérique du noyau.
\end{itemize}

La chaîne exécutable est le répertoire
\texttt{../python/scha\_pipeline/}, distinct des prototypes archivés. Son point
d'entrée de validation, depuis la racine de la livraison, est
\begin{verbatim}
python3 python/scha_pipeline/validate_pipeline.py --quick
\end{verbatim}
La même commande sans \texttt{--quick} exécute également les contrôles de
convergence les plus coûteux. Elle est non destructive : les caches et fichiers
temporaires sont redirigés, et tout échec obligatoire produit un code de retour
non nul. Le fichier \texttt{python/scha\_pipeline/requirements.txt} fixe les
versions des bibliothèques scientifiques utilisées avec Python 3.9.6 ; le fichier
\texttt{python/scha\_pipeline/README.md} donne le graphe de calcul, les unités,
les scripts appelés et un exemple de commande de bout en bout.

Un premier calcul autonome est lancé par
\begin{verbatim}
python3 python/scha_pipeline/run_pipeline.py
\end{verbatim}
Il lit \texttt{inputs/examples/batio3\_500K.json} et écrit la solution et la
réponse dans \texttt{outputs/example/}. Les nouveaux calculs sont séparés des
sorties de référence ; le README décrit également les scans utilisés pour les
figures du rapport et de la soutenance.
